% Principal: GDZ PPN236018647, 1917, PDF native 13, printed 6.
\SourceLeaf{13}\NoSourceHead
\pdfbookmark[1]{Eine in Deutschland erfundene Rechenmaschine}{rechenmaschine}
\HistoricTitle{EINE IN DEUTSCHLAND ERFUNDENE RECHEN-\\MASCHINE.}
\HistoricCitation{\textsc{Dinglers} \textit{Polytechnisches Journal} Bd. 52 (1834), S. 237,\\
entnommen der Zeitschrift \textit{Le National}, 27. März 1834.}

Herr \textsc{Schiereck}, Professor der Mathematik zu Frankfurt a. M., hat der französischen Akademie der Wissenschaften eine Dissertation über die Theorie der Zahlen eingeschickt. Dieser Abhandlung ist ein Zeugnis des Herrn \textsc{Gauss}, des berühmten Geometers zu Göttingen, beigelegt, folgenden Inhalts:

Herr \textsc{Schiereck} hat mir ein Modell einer Rechenmaschine gezeigt, welche er zur Ausführung der arithmetischen Operationen erfunden hat. Ich bezeuge mit Vergnügen, dass diese Maschine den beabsichtigten Zweck sehr leicht erreicht, und dass dieses nach den Verbesserungen, welche der Erfinder an ihr zu machen beabsichtigt, noch mehr der Fall sein wird. Diese sinnreiche Erfindung ist umso schätzbarer, weil diese Maschine mit geringen Kosten hergestellt werden kann.

\bigskip\begin{center}\small\textbf{BEMERKUNG.}\end{center}
\begingroup\small\par\indent
Die vorstehende Mitteilung hat R. \textsc{Mehmke} im 48. Bande der Zeitschrift für Mathematik und Physik S. 318 wieder abdrucken lassen. Joseph Friedrich \textsc{Schiereck} aus Posen gebürtig, hatte in Göttingen studiert. Über seine Rechenmaschine hat sich nichts näheres feststellen lassen. Er starb im November 1842. Vergl. H. E. \textsc{Scriba}, Biographisch-literarisches Lexikon der Schriftsteller des Großherzogtums Hessen im 19. Jahrhundert; 2. Abteilung, S. 638 ff.
\begin{flushright}\textsc{Schlesinger}.\end{flushright}
\endgroup
\ClosingRule
\clearpage
